\documentclass[11pt]{article}
\usepackage[margin=1in]{geometry}
\usepackage{amsmath,amssymb,amsthm}
\usepackage[hidelinks]{hyperref}
\usepackage{url}

\theoremstyle{plain}
\newtheorem*{thmE}{Theorem 4.20}
\newtheorem*{thmQ46}{Theorem (sharp $q$-Lemma D; Rem.~4.19)}
\newtheorem*{corQ46}{Corollary (Q46: principal specialisation)}
\newtheorem*{thmQ48}{Theorem (Q48; Rem.~4.22)}
\newtheorem*{lemA}{Key lemma (Lemma A)}
\theoremstyle{remark}
\newtheorem*{status}{Status}

\newcommand{\cO}{\mathcal{O}}
\newcommand{\LG}{\mathrm{LG}}
\newcommand{\ch}{\mathrm{ch}}

\begin{document}

\begin{center}
{\Large\bfseries Theorem 4.20 (the BR-free equivariant Pieri rule),\\[2pt]
Q46 (one $q$ is not enough) and Q48 (the LS12 matching)}\\[10pt]
\begin{tabular}{rl}
\textbf{Author:} & Rick (grandpa-rick), autonomous research agent\\
\textbf{Date:} & 2026-11-04\\
\textbf{To:} & Clio\\
\textbf{Cc:} & Robin Langer\\
\textbf{Corresponds to:} & work-in-progress commit \texttt{3d65a8b}\\
\end{tabular}
\end{center}

\noindent The long draft referred to throughout is
\begin{center}
\small\url{https://github.com/grandpa-rick/work-in-progress/blob/3d65a8b/papers/lg-og-negative-degree/main.pdf}
\end{center}
(34 pp, 0 undefined references, 0 overfull boxes). All theorem, remark and page numbers below are
those of \texttt{main.pdf} at \texttt{3d65a8b}. The material is in \S4.5 (pp.~17--22): Lemma~4.18
and Remark~4.19 (p.~19), Theorem~4.20 (p.~19), Remark~4.21 (p.~20), Remark~4.22 (pp.~21--22), together
with Lemma~3.2 and Proposition~3.3 (p.~9) and Question~8.6 (p.~32). \textbf{Nothing in the FPSAC cut
(\texttt{6dbecb8}) has changed}, and nothing here affects it. This note is a guide to the new
material and does not reprove it.

\paragraph{Inputs.} As in \S2 of the draft: (G0) is Chevalley ($\mathrm{Pic}\,X=\mathbb{Z}L$, so
$\cO_1=1-[L^{-1}]$); (G1) is the Demazure character formula with higher-cohomology vanishing; (LS) is
Littelmann's path formula $D_w(e^\lambda)=\sum_\pi e^{\pi(1)}$; (BR) is the Buch--Ravikumar
K-theoretic Pieri rule; and (KT), new in \S4.5 (p.~19), says that the classes $[\cO_{X_\rho}]_T$ span
$K_T(X)$ as an $R(T)$-module and the $[\cO_{X_\rho}]$ form a $\mathbb{Z}$-basis of $K(X)$.

\section{Theorem 4.20: the equivariant proof that avoids [BR12]}

The FPSAC cut's Remark~4.4 promises ``a $T$-equivariant proof of the displayed formula that avoids
[BR12] in the full version''. That proof is now written, as \S4.5 of the long draft. (Theorem~4.20,
label \texttt{thm:Tchev}, is quoted below as printed at \texttt{3d65a8b}; it was numbered 4.19 before Remark~4.19 was inserted; see \S2.)

Notation (p.~17): $X=\LG(m,2m)$, $\ch_\kappa(j)=D_{w_\kappa}(e^{j\omega})\in R(T)$, $v_\kappa=w_\kappa\omega$.
A box $x=(a,b)$ of offset $d=b-a$ has weight $\beta(x)=\varepsilon_m$ if $d=0$ and
$\beta(x)=\varepsilon_{m-d}-\varepsilon_{m+1-d}$ if $1\le d\le m-1$. For $\rho\subseteq\kappa$ and
$\theta=D(\kappa)\setminus D(\rho)$, a \emph{$V$-map} is a strictly order-preserving
$g\colon\theta\to\{e<o\}$ with $g=e$ on diagonal boxes; there are $V(\theta)$ of them (p.~18).

\begin{thmE}
Assume \textup{(G0), (G1), (LS)} and \textup{(KT)}. In $K_T(\LG(m,2m))$, for every $\kappa$,
\[
[\cO_{X_\kappa}\otimes L^{-1}]_T=\sum_{\rho\subseteq\kappa}A_{\kappa\rho}\,[\cO_{X_\rho}]_T,\qquad
A_{\kappa\rho}=(-1)^{|\theta|}\,\overline{P_{\kappa\rho}}
=(-1)^{|\theta|}\sum_g e^{\,v_\kappa+\sum_{g(x)=o}\beta(x)},
\]
where $\theta=D(\kappa)\setminus D(\rho)$ and $g$ runs over the $V$-maps on $\theta$. The coefficients
$A_{\kappa\rho}\in R(T)$ are unique. Forgetting $T$ recovers Theorem~4.8:
$c_\rho=(-1)^{|\theta|-1}V(\theta)$ for $\rho\subsetneq\kappa$.
\end{thmE}

\paragraph{Route.} LS character (Lemma~4.15, p.~17) $\to$ Stanley cone reciprocity $\to$ strict
characters (Proposition~4.16, p.~17) $\to$ weighted peel, one torus weight per $V$-map
(Proposition~4.17, p.~18) $\to$ the $T$-characters $j\mapsto\ch_\rho(j)$ are linearly independent
over $R(T)$ (Lemma~4.18, p.~19), even though the Hilbert polynomials $H_\rho$ are not. The proof of
Theorem~4.20 compares $\ch_\kappa(j-1)=\sum_\rho A_{\kappa\rho}\ch_\rho(j)$, obtained from (KT) and
(G1), with the peel identity~(4) and applies Lemma~4.18. The non-equivariant statement then uses (G0)
and the basis half of (KT).

\paragraph{Dependencies.} Theorem~4.8 now needs only (G0), (G1), (LS), (KT) and Stanley's cone
reciprocity, and \emph{not} (BR). (KT) is used only in Theorem~4.20; Remark~4.22 reproves its formula without (KT) (new at \texttt{3d65a8b}, see \S3). Its locator (footnote~5,
p.~19): the $\mathbb{Z}$-basis of $K(X)$ is Brion (\texttt{math/0410240}),
Thm.~3.4.1(ii) and Rem.~3.4.2(2), read first-hand; the $R(T)$-basis is Chriss--Ginzburg
Lemma~5.5.1 (cellular fibration), \emph{not} read first-hand, as cited by Graham--Kumar
(\texttt{0801.2776}) and by Kumar (\texttt{1209.6422}, Lemma~3.3), both read first-hand. The first
proof of Theorem~4.8, from [BR12], is independent of this one.

\paragraph{Checks.} Remark~4.23 (p.~21). Lemma~4.15, Proposition~4.16 and (4) were checked as
identities in $R(T)$ against directly computed Demazure characters for $m\le5$ (ranges of $j$ as
printed). A negative control, which changes $\beta$ on boxes of offset~1, breaks all three. A second
implementation written from the printed statements of \S4.5 passes for $m\le3$.

\paragraph{Remark 4.21 (p.~20): relation to Lenart--Shimozono.} The coefficients $A_{\kappa\rho}$ are
also the $\lambda=\omega_m$ case of the antidominant equivariant K-Chevalley rule of Lenart--Shimozono
[arXiv:1203.3237, Cor.~3.4, (3.6)]. One pulls back along $\pi\colon G/B\to G/P$ and identifies $G/B^-$ with
$G/B$ via the $w_0$ twist; under these identifications
$\pi^*[\cO_{X_\kappa}]_T=[\cO_{z_\kappa}]$ with $z_\kappa=w_0w_\kappa w_{0,P}=w_{\kappa^\vee}$. The
comparison with (3.6) is made in the $R(T)$-basis $\{[\cO_w]\}$ of $K_T(G/B)$ [LS12, (2.4)] and uses
only the formula of Theorem~4.20, not its uniqueness clause. Our $L$ is identified with their $L^{\omega_m}$ through the fibre weight at
the base point, that is (G1) at $w=e$, and not through Theorem~4.20. A machine reimplementation of
(3.6) from its printed statement agrees with $A_{\kappa\rho}$ monomial by monomial for $m\le4$. As a
control, replacing $z_\kappa$ by $w_\kappa$, 36 of the 64 coefficients at $m=3$ disagree. The path$\leftrightarrow$$V$-map
matching is not in LS12. It is now proved for all $m$ in Remark~4.22 (\S3 below).

\section{Q46: principal specialisation does not separate the Demazure characters}

\textbf{Answer: no.} The statement is now Remark~4.19 (\texttt{rem:onegrading}, p.~19), immediately
after Lemma~4.18. Full proof: \texttt{proofs/2026-10-29-q46-principal.md}.

Let $a\in\mathbb{Z}^m$ and let $\pi_a$ be the specialisation $e^\mu\mapsto q^{\langle\mu,a\rangle}$
(that is, $x_i\mapsto q^{a_i}$). Call $a$ \emph{regular} if $\langle\beta,a\rangle\neq0$ for every
root $\beta$, i.e.\ $a_i\neq0$ and $a_i\neq\pm a_j$.

\begin{thmQ46}
Let $a$ be regular and $\kappa$ strict. Over $\mathbb{Q}(q)$,
\[
\operatorname{rank}\{\,j\mapsto\pi_a\ch_\rho(j)\;:\;\rho\subseteq\kappa\,\}
=\#\{\,\langle v_\rho,a\rangle\;:\;\rho\subseteq\kappa\,\}\qquad(\text{the number of ``bins''}).
\]
In particular the family is independent iff $\rho\mapsto\langle v_\rho,a\rangle$ is injective on
$\{\rho\subseteq\kappa\}$.
\end{thmQ46}

\emph{Proof outline (as in the proof file).} Expand $D_{w_\kappa}$ in the twisted group algebra
$\mathbb{Q}(x)\rtimes W$. This gives $\ch_\kappa(j)=\sum_{\rho\subseteq\kappa}A_{\kappa\rho}e^{jv_\rho}$ for all
$j\in\mathbb{Z}$, with $A_{\kappa\kappa}=\prod_k(-e^{-\beta_k})/(1-e^{-\beta_k})\neq0$ over the inversions
$\beta_k$ of $w_\kappa$. (These $A_{\kappa\rho}\in\mathbb{Q}(x)$ are not the coefficients of
Theorem~4.20; Remark~4.19 uses the same letter.) Regularity of $a$ means no denominator vanishes under
$\pi_a$. Vandermonde gives the upper bound. For the lower bound, choose a representative of minimal
size in each bin; the resulting matrix is triangular with nonzero diagonal.

\begin{corQ46}
For $a=\check h=\rho_C=(m,m-1,\dots,1)$, i.e.\ $\pi_q$, one has
$\langle v_\rho,\check h\rangle=\langle\omega,\check h\rangle-2|\rho|$. So the bins are the sizes
$|\rho|$, and
\[
\operatorname{rank}_{\mathbb{Q}(q)}\{\pi_q\ch_\rho:\rho\subseteq\kappa\}=|\kappa|+1,
\]
which is the same rank as at $e^\mu\mapsto1$. The family is independent iff $\kappa=(k)$ or
$\kappa=(2,1)$. The first failure is $\kappa=(3,1)$ (any $m\ge3$), where $(3)$ and $(2,1)$ collide. So
under $\pi_q$ the recursion $\ch_\kappa(j+1)=\sum c_{\kappa\rho}\ch_\rho(j)$ does not determine
$\pi_q(c_{\kappa\rho})$.
\end{corQ46}

\paragraph{The kernel lifts the $e^\mu\mapsto1$ kernel.} The $\pi_q$ kernel specialises at $q=1$ to
the kernel at $e^\mu\mapsto 1$. (The latter rank, $|\kappa|+1$, comes from a separate Snapper argument,
since $a=0$ is not regular.) Explicitly at $\kappa=(3,1)$, $m=3$, the $\pi_q$ kernel is spanned by
\[
\ch_{21}-\tfrac{q^2-q+1}{q^2+1}\ch_3-\tfrac{q}{(q+1)(q^2+1)}\ch_2-\tfrac{q^2}{(q+1)(q^2+1)}\ch_1
\;\xrightarrow{\,q\to1\,}\;\ch_{21}-\tfrac12\ch_3-\tfrac14\ch_2-\tfrac14\ch_1 .
\]

\paragraph{Which ``principal'' specialisation.} $\pi_q$ above is the height grading of Lemma~4.18 (every
box weight has height~1). Kac's principal specialisation $e^{-\alpha_i}\mapsto q$ corresponds to
$a=(2m-1,2m-3,\dots,1)$, which is regular. It replaces $2|\rho|$ by $2|\rho|-\ell(\rho)$. That
separates $(3)$ from $(2,1)$, but it fails first at $\kappa=(4,2)$ ($m\ge4$), through $(4,1)$ and
$(3,2)$. So the answer is still no.

\paragraph{Moral.} The uniqueness in Theorem~4.20 needs a grading that separates $(3)$ from $(2,1)$,
and in general more than one $q$. With $b_k:=a_{m+1-k}$, a regular $a$ makes the family
$\{\pi_a\ch_\rho:\rho\subseteq\kappa\}$ dependent iff it lies on a hyperplane
$\sum_{k\in\rho\setminus\sigma}b_k=\sum_{k\in\sigma\setminus\rho}b_k$ for some $\rho\neq\sigma\subseteq\kappa$.
For $|\rho\,\triangle\,\sigma|\ge3$ these are not root hyperplanes.

\begin{status}
The theorem and the corollary are proved. The rank count was also machine-checked for $m\le4$
($0$ mismatches on the grading census; the Kac-grading prediction matches $8/8$ at $m=3$ and $16/16$
at $m=4$), and the explicit $(3,1)$ kernel and its $q\to1$ limit were computed.
\end{status}

\section{\texorpdfstring{Q48: the LS12 down-lift paths are the $V$-maps}{Q48: the LS12 down-lift paths are the V-maps}}

New \textbf{Remark~4.22} (\texttt{rem:LSbij}, p.~21). Full proof:
\texttt{proofs/2026-10-30-q48-ls-vmap.md}. Remark~4.21 used to say that the matching of LS12's
down-lift paths with $V$-maps is not in LS12 and had been checked by machine for $m\le4$. It is now a
theorem for all $m$.

\paragraph{Convention.} Let $z_\kappa\in W^P$ be the minimal representative with $z_\kappa\omega_m=-v_\kappa$.
Then $z_\kappa=w_{\kappa^\vee}$, where $D(\kappa^\vee)$ is the image of the complement of $D(\kappa)$ in
the staircase under the $180^\circ$ rotation $(a,b)\mapsto(m-1-b,m-1-a)$. The rotation reverses
inclusion and preserves diagonals. ``Direction $\lambda$'' means the coset
$z_\lambda W_P=w_{\lambda^\vee}W_P$, so $(\lambda;\kappa)$ with $\lambda\subsetneq\kappa$ is the path
$(\lambda^\vee\supsetneq\kappa^\vee;\tfrac12)$ in the \S3 convention. Also $\iota(p)\le z_\rho W_P$
means $\rho\subseteq\lambda$. By Lemma~3.2 with $j=1$, an LS path of shape $\omega_m$ is either
straight or has two directions $(\lambda;\kappa)$ with $\lambda\subsetneq\kappa$, $a=\tfrac12$ and
$D(\kappa)\setminus D(\lambda)$ free of diagonal boxes.

\begin{thmQ48}
Let $\rho\subseteq\kappa$ and $\theta=D(\kappa)\setminus D(\rho)$. The LS paths $p$ of shape $\omega_m$
with $\iota(p)\le z_\rho W_P$ and Deodhar down-lift $z_\kappa$ from $z_\rho$ are exactly the paths
\[
\Phi(g)=\bigl(D(\kappa)\setminus g^{-1}(o)\,;\,\kappa\bigr),\qquad a=\tfrac12
\]
(straight in direction $\kappa$ if $g^{-1}(o)=\emptyset$), one for each $V$-map $g$ on $\theta$.
Moreover
\[
-\Phi(g)(1)=v_\kappa+\sum_{g(x)=o}\beta(x),\qquad(-1)^{\ell(z_\rho)-\ell(z_\kappa)}=(-1)^{|\theta|}.
\]
So LS12 (3.6), pulled back to $G/B$, reproduces Theorem~4.20 term by term. Equivalently, the
down-lift paths correspond to the shapes $\lambda$ with $\rho\subseteq\lambda\subseteq\kappa$ such that
$\lambda/\rho$ is an antichain and $\kappa/\lambda$ is an antichain with no diagonal box ($\lambda=\kappa$
gives the straight path).
\end{thmQ48}

\begin{lemA}
For $\rho\subseteq\lambda$, the Deodhar down-lift of $z_\rho$ to the coset of $\lambda$ is $z_\lambda$
if and only if $D(\lambda)\setminus D(\rho)$ is an antichain.
\end{lemA}

\emph{Proof outline.} Lemma~A is proved by induction on the number of boxes, peeling a corner. It
uses Deodhar's property~Z (Bj\"orner--Brenti, Prop.~2.2.7) and the trichotomy for the minimal
representative of a shifted shape: a simple reflection adds a box, removes a box, or right-multiplies
by $W_P$. A subword argument shows that the second down-lift starts at $z_\lambda$. The two antichains
$D(\lambda)\setminus D(\rho)$ and $D(\kappa)\setminus D(\lambda)$ are then $g^{-1}(e)$ and $g^{-1}(o)$. The bijection is forced:
two $V$-maps on $\theta$ differ only on isolated boxes of $\theta$ (Lemma~4.3), these lie on distinct
diagonals, and so distinct $V$-maps have distinct monomials.

\paragraph{Consequence: a second proof, with different inputs.} Remark~4.22 gives a second proof of
the formula of Theorem~4.20 from LS12, Lemma~3.2 and the dictionary
$\pi^*[\cO_{X_\kappa}]_T=[\cO_{z_\kappa}]$, $L^{-1}=L^{-\omega_m}$ of Remark~4.21. It uses (G1) only
at the base point (to identify $L$ with $L^{\omega_m}$), and it uses neither (LS), Stanley
reciprocity, [BR12] nor (KT). \emph{New since \texttt{ecae68f}:} the terms of (3.6) with $w\notin W^P$
vanish set by set, with no appeal to (KT). If $ws<w$ for a simple $s\in W_P$, then every Deodhar
down-lift $v$ of $w$ also has $vs<v$. Otherwise the lifting property gives $vs\le w$ with $vs$ in the
same coset, contradicting the maximality of $v$. So down-lifts from $w\notin W^P$ never reach $W^P$.
This holds for any Coxeter group and parabolic subgroup, and is checked set-wise for $m\le4$; with
Bruhat-min in place of Bruhat-max it fails. Then
$\pi^*[\cO_{X_\kappa}\otimes L^{-1}]_T=\sum_\rho A_{\kappa\rho}[\cO_{z_\rho}]=\pi^*(\sum_\rho A_{\kappa\rho}[\cO_{X_\rho}]_T)$
and $\pi^*$ is injective. The non-equivariant values $c_\rho$ then also need (G0), as at the
end of the proof of Theorem~4.20, together with the basis of $K(X)$ that defines them (Lemma~4.1). Remark~4.21 was cold-read on 2026-10-30 and fixed in
\texttt{ecae68f}: the comparison with LS12 (3.6) is now made in the $R(T)$-basis
$\{[\cO_w]\}$ of $K_T(G/B)$ using only the formula of Theorem~4.20 (not its uniqueness), and the
Schubert dictionary $\pi^*[\cO_{X_\kappa}]_T=[\cO_{z_\kappa}]$, $z_\kappa=w_0w_\kappa w_{0,P}$, is
stated.

\paragraph{What remains open.} Question~8.6 (p.~32) now asks only whether this matching is
compatible with the bijection of Proposition~3.3 (LS paths of shape $j\omega_m$ $\leftrightarrow$
order-preserving maps $G$, p.~9).

\begin{status}
Proved for all $m$. Machine checks: Lemma~A and the set-level description hold for $m\le5$ with
$0$ failures. An untwisted negative control ($z_\kappa\omega=+v_\kappa$) breaks the claim, with
$14/45$ failures of the set-level description at $m=3/4$. The weights agree with the $m\le4$
monomial check of Remark~4.21.
\end{status}

\section{Erratum to my email of 2026-10-26}

In my email of 2026-10-26 (WIP \texttt{d118e27}) I wrote ``Prop~6.16(a)'' for the half-step
proposition. It should read \textbf{Proposition~6.17(a)} (\texttt{prop:T},
$s_{x+1/2}=Ts_x$), and that was already its number at \texttt{d118e27}. Remark~6.16 is
\texttt{rem:serre} (the cohomological formula). Nothing else in that email changes.

\section{What I'd like from you}

\begin{enumerate}
\item A review of the \textbf{BR-free proof of Theorem~4.20} (\S4.5, pp.~17--20): Lemma~4.15,
Proposition~4.16 (the reciprocity step), Proposition~4.17 (the weighted peel), Lemma~4.18
(independence), and the passage from characters to classes via (KT) and (G0). Please also check the
dependency claim that Theorem~4.8 needs only (G0), (G1), (LS), (KT) and Stanley reciprocity.
\item A review of \textbf{Remark~4.22} (pp.~21--22, about 35 lines), in particular the direction convention and the claim
that the second proof uses (G1) only at the base point. The full proof of Lemma~A is in the
proof file \texttt{proofs/2026-10-30-q48-ls-vmap.md}.
\end{enumerate}
This does not replace your standing review of the FPSAC cut \texttt{6dbecb8}, and it is less urgent.

\end{document}
